% =====================================================================
% 12d_glossaire.tex — glossaire des termes techniques
% =====================================================================
\chapter{Glossaire}
\label{ann:glossaire}

\section{Termes du modèle objet}

\begin{description}[leftmargin=2.6em,style=nextline]

\item[Attribut objet.]
Donnée membre d'un type objet. Dans le projet, \code{client\_type} possède
notamment \code{nom}, \code{adresse} et \code{ref\_parrain}.

\item[Collection.]
Ensemble ordonné ou non d'éléments de même type. Le projet utilise un
\code{VARRAY} pour les téléphones et les mois, et une \emph{nested table} pour
les signataires et les statistiques.

\item[Héritage.]
Mécanisme par lequel un sous-type reprend les attributs et méthodes d'un
super-type. \code{cptCourant\_type}, \code{cptEpargne\_type} et
\code{cptTerme\_type} héritent de \code{compte\_type}.

\item[Méthode membre.]
Fonction ou procédure associée à une instance de type. Elle peut lire
\code{SELF}. \code{passeMvt} est une méthode membre de
\code{cptCourant\_type}.

\item[Méthode statique.]
Fonction ou procédure associée au type, et non à une instance. Elle est appelée
avec le nom du type. \code{fermeCompte} est une méthode statique de
\code{compte\_type}.

\item[Nested table.]
Collection non ordonnée, stockée dans une table de stockage. Les signataires
d'un compte courant et les statistiques d'un mois sont des \emph{nested
tables}.

\item[Objet.]
Instance d'un type objet, avec ses attributs et ses méthodes. Une ligne de
\code{cpt\_courant} est un objet \code{cptCourant\_type}.

\item[OID (identifiant objet).]
Identifiant interne ou dérivé d'une clé qui identifie un objet persistant. La
table \code{client} utilise un OID système ; la table \code{agence} utilise
\code{OBJECT IDENTIFIER IS PRIMARY KEY}.

\item[Polymorphisme.]
Capacité à appeler une méthode dont l'implémentation dépend du type réel de
l'objet. \code{fraisTenue} est redéfinie dans chaque sous-type de compte.

\item[REF.]
Référence persistante vers un objet. Elle se déréférence par la notation
pointée, par exemple \code{c.ref\_conseiller.nom}.

\item[Type objet.]
Définition d'une structure de données avec attributs et méthodes. Le projet
distingue les types métier, les types de collection et les types de vue.

\item[VARRAY.]
Collection ordonnée de taille maximale fixe. \code{telephones} est un
\code{VARRAY(3)}.

\end{description}

\section{Termes relationnels et SQL}

\begin{description}[leftmargin=2.6em,style=nextline]

\item[Contrainte différée.]
Contrainte dont la vérification est reportée au \code{COMMIT}. La contrainte
\code{ck\allowbreak\_cc\allowbreak\_client} est
\code{DEFERRABLE INITIALLY DEFERRED}.

\item[FK (clé étrangère).]
Contrainte qui garantit qu'une valeur référence une clé existante. Elle peut
porter sur une colonne ou sur un attribut \code{REF}.

\item[Index.]
Structure d'accès qui accélère les recherches. Les clés primaires et les
contraintes d'unicité créent des index.

\item[PDB (pluggable database).]
Base de données pluggable d'Oracle. Le projet utilise \code{XEPDB1}.

\item[Séquence.]
Objet qui produit des numéros successifs. Le projet utilise
\code{seq\_operation}, \code{seq\_alerte}, \code{seq\_audit} et
\code{seq\_affectation}.

\item[SQL*Plus.]
Client en ligne de commande d'Oracle. Les scripts utilisent ses paramètres
\code{SERVEROUTPUT}, \code{LINESIZE} et \code{DEFINE OFF}.

\item[Vue objet.]
Vue dont les lignes sont des objets d'un type donné. \code{Agence\_VOR} est une
vue objet avec un OID fondé sur le code d'agence.

\end{description}

\section{Termes PL/SQL et transactions}

\begin{description}[leftmargin=2.6em,style=nextline]

\item[Bloc anonyme.]
Bloc PL/SQL sans nom, exécuté directement. Les tests D7 et D9 utilisent des
blocs anonymes.

\item[COMMIT.]
Validation définitive de la transaction courante. \code{fermeCompte} et
\code{drop\_client} valident eux-mêmes ; \code{virement} et \code{passeMvt} ne
le font pas.

\item[Curseur.]
Pointeur sur le résultat d'une requête. Les blocs D7 utilisent des curseurs
explicites et paramétrés.

\item[Exception.]
Erreur détectée par PL/SQL ou levée par \code{RAISE\_APPLICATION\_ERROR}. Les
codes \code{ORA-20001} à \code{ORA-20030} sont applicatifs.

\item[ROLLBACK.]
Annulation de la transaction courante. Les tests de rollback prouvent que
\code{virement} et \code{passeMvt} ne valident pas implicitement.

\item[SAVEPOINT.]
Point de restauration intermédiaire dans une transaction. \code{virement}
l'utilise pour annuler les deux écritures en cas d'échec.

\item[Trigger.]
Programme stocké déclenché par un événement sur une table ou une vue. Le projet
compte six triggers.

\item[Trigger COMPOUND.]
Trigger qui combine plusieurs sections temporelles :
\code{BEFORE STATEMENT}, \code{BEFORE EACH ROW}, \code{AFTER EACH ROW} et
\code{AFTER STATEMENT}. Il évite l'erreur de table mutante
\code{ORA-04091}.

\end{description}

\section{Termes du projet}

\begin{description}[leftmargin=2.6em,style=nextline]

\item[Agence.]
Structure qui emploie des conseillers et détient des comptes. Le projet en
compte trois : A1, A2 et A3.

\item[Alerte.]
Message stocké dans la table \code{alerte}, par exemple
\code{SOLDE\_NEGATIF}, \code{OPERATION\allowbreak\_SUSPECTE} ou
\code{ECHEANCE\allowbreak\_PROCHE}.

\item[Audit.]
Trace des modifications sensibles. La table \code{"AUDIT"} est entre guillemets
car \code{AUDIT} est un mot réservé.

\item[Conseiller.]
Employé qui suit des clients et peut être affecté à une agence. Le projet en
compte quatre.

\item[Compte courant.]
Compte utilisé pour les dépôts, retraits, CB et virements. Il possède un
découvert autorisé et des signataires.

\item[Compte d'épargne.]
Compte rémunéré avec un taux d'intérêt et un plafond de dépôt.

\item[Compte à terme.]
Compte dont le montant est bloqué jusqu'à une date d'échéance.

\item[Signataire.]
Client autorisé à opérer sur un compte courant, avec un droit \code{X},
\code{R} ou \code{D}.

\item[Virement.]
Transfert entre deux comptes courants, avec commission inter-agences de 1~\%,
minimum 100~DA.

\end{description}
